\chapter{A preparation and randomized-readout assay}
\label{ch:rr-preparation}

The previous finite-pulse result used a stationary preparation and a selective
change to one physical transition channel.  Its positive first-hitting kernel
also supports a different experiment: prepare a detected state, compare it with
another preparation, and randomize the time of a single endpoint reading.  This
removes the stationary-preparation and selective-control requirements.  It
replaces them with a trusted reset and a common, time-homogeneous evolution.
Neither access contract is automatically easier to establish in an apparatus.

The mathematical result below is a Markov occupation maximum principle.
Converting it into an operational assay is useful, but it does not establish the
first test of arbitrary-size classical dynamics.  Established survival
inequalities already address closely related questions
\cite{rrLambert2010,rrLi2012,rrZhou2015}.  Section~\ref{sec:rr-comparison} compares
the present assay with those alternatives and records a baseline that outperforms
its simplest version.

\section{The classical preparation inequality}
\label{sec:rr-preparation-theorem}

Let $Q$ be a finite, time-homogeneous row generator and $P(t)=e^{tQ}$.
A reference arm prepares state $j$ exactly.  A comparison arm prepares a fixed
probability distribution $\rho$ over the same states.  Both arms evolve under
$Q$.  Independently of the trajectory and preparation arm, draw a readout time
uniformly from $[0,T]$, measured from preparation.  Record whether the system
occupies the singleton state $j$.

\begin{theorem}[Integrated return maximum]
\label{thm:rr-return}
For $i\ne j$ and $T>0$,
\begin{align}
 K_{ij}(T)
 &=\int_0^T\bigl(P_{jj}(t)-P_{ij}(t)\bigr)\,dt\nonumber\\
 &=\int_0^T P_{jj}(s)\Pr_i(\tau_j>T-s)\,ds>0,
 \label{eq:rr-hitting}
\end{align}
where $\tau_j$ is the first hitting time of $j$ from $i$.
Consequently, for every initial distribution $\rho$,
\begin{equation}
 C_{\rho,j}(T)
 =\frac1T\int_0^T\left(P_{jj}(t)-[\rho P(t)]_j\right)\,dt
 \ge0.
 \label{eq:rr-classical-null}
\end{equation}
The inequality is strict if $\rho_j<1$.  Irreducibility and a stationary
preparation are unnecessary.  If the exit rates from $i,j$ are at most
$\Lambda$, then $K_{ij}(T)/T\ge e^{-\Lambda T}$.
\end{theorem}

\begin{proof}
The strong Markov property at $\tau_j$ gives
\[
 P_{ij}(t)=\int_{[0,t]}P_{jj}(t-u)\Pr_i(\tau_j\in du).
\]
Integrating this identity and interchanging nonnegative integrals gives
\eqref{eq:rr-hitting}.  No departure from $j$ for $s$ time units has probability
at least $e^{-\Lambda s}$, and no departure from $i$ for $T-s$ units has
probability at least $e^{-\Lambda(T-s)}$.  These events give the claimed lower
bound and strict positivity.  Finally,
$T C_{\rho,j}=\sum_i\rho_iK_{ij}$, with $K_{jj}=0$.
\end{proof}

Only the reference preparation must be a trusted singleton.  The comparison
preparation may be an unknown fixed distribution.  This distinction matters for
calibration: uncertainty about that comparison distribution need not enlarge
the classical null, although it affects predicted power against a specified
alternative.  A computational-basis label in a quantum experiment does not by
itself establish a singleton classical hidden state.

A sufficient coarse-target extension is a genuine reset interface.  Let $B$ be
the detected set, with reference preparation $\nu_B$ supported in $B$ and a
comparison preparation outside $B$.  Suppose the conditional state at first
entry into $B$ is $\nu_B$, independently of the entry time and previous history.
The strong Markov argument then replaces $P_{jj}(s)$ by
$\Pr_{\nu_B}(X_s\in B)$ and proves the same inequality.  Mixtures of an outside
preparation with $\nu_B$ are also allowed.  Arbitrary states initially inside
$B$ are not covered by this sufficient condition.

\section{Protected clocks and failure controls}
\label{sec:rr-clocks}

\begin{theorem}[Decreasing readout densities]
\label{thm:rr-decreasing}
Let $w$ be a nonnegative, normalized, nonincreasing readout-time density on
$[0,\infty)$, shared by both preparations and independent of the dynamics.
Under Theorem~\ref{thm:rr-return},
\[
 \int_0^\infty w(t)
 \left(P_{jj}(t)-[\rho P(t)]_j\right)\,dt\ge0.
\]
Every such density is, up to null sets, a mixture of uniform densities:
\begin{equation}
 w(t)=\int_{[t,\infty)}\frac{\mu(dT)}{T},
 \qquad \mu(dT)=-T\,dw(T),
 \label{eq:rr-mixture}
\end{equation}
where $\mu$ is a probability measure.
\end{theorem}

\begin{proof}
For a decreasing integrable density, $t w(t)\to0$ at both endpoints.
Stieltjes integration by parts therefore gives
$\int\mu=\int_0^\infty w(t)\,dt=1$ and
$\int_{[t,\infty)}\mu(dT)/T=w(t)$, with endpoint conventions irrelevant to
probability densities.  Tonelli's theorem expresses the asserted contrast as
$\int C_{\rho,j}(T)\mu(dT)$, which is nonnegative by
Theorem~\ref{thm:rr-return}.  Equivalently, integration by parts against the
nonnegative cumulative kernel $K$ proves the sign directly.
\end{proof}

The mixture representation is established mathematics
\cite{rrPavlides2012}; its role here is to expose experimental timing choices.
Randomizing a cutoff $T$ and then drawing time uniformly in $[0,T]$ preserves
the inequality exactly.  For a specified alternative, no decreasing-density
clock can beat the best uniform horizon unless an additional constraint changes
the optimization problem.  Mean-time constraints, for example, lead to a
standard convex design problem.

The following controls show why the assumptions cannot be omitted.

\begin{enumerate}
\item \emph{A fixed or delayed readout can reverse the classical contrast.}
A directed three-cycle has $P_{jj}(t)<P_{ij}(t)$ at some times.  Its uniform
average over a delayed interval can also be negative.  Uniform sampling must
start at preparation.  Increasing time densities have no general guarantee.
\item \emph{A coarse detector can reverse the contrast.}
Take four states $j,i,k,a$, with rates $j\to a$ and $i\to k$ both equal to $2$,
and $a,k$ absorbing.  Detect $B=\{j,k\}$, prepare $j$ or $i$, and average over
$[0,1]$.  The contrast is
\[
 \int_0^1\bigl(2e^{-2t}-1\bigr)\,dt=-e^{-2}<0.
\]
This wholly classical model violates the singleton conclusion because entry
into $B$ does not reset to the reference state.
\item \emph{Scheduled dynamics can reverse the contrast.}
During $[0,1]$ allow only $j\to a$ at rate $4$; during $[1,2]$ allow only
$i\to j$ at rate $4$.  Preparing $j$ or $i$ and averaging the singleton-$j$
readout over $[0,2]$ gives
\[
 \frac12\left[\frac{2(1-e^{-4})}{4}+e^{-4}-1\right]<0.
\]
The generator is time dependent.  Augmenting the system with an internal clock
does not restore the required singleton or reset observation contract.
\end{enumerate}

Nonincreasing density is a sufficient condition, not a necessary
characterization of all safe clocks.  For example, placing at least half of the
readout probability at time zero protects the sign against any remaining
stochastic transition kernel.  The general clock cone is not characterized by
this work.  A rejection of the declared classical contract is not a
device-independent identification of quantum dynamics.

\section{A damped quantum model and an exact threshold}
\label{sec:rr-quantum}

Consider the valid qubit Lindblad generator
\begin{equation}
 \mathcal L(\rho)=-i\frac{\omega}{2}[X,\rho]
  +\gamma\left(\operatorname{Tr}(\rho)\frac I2-\rho\right),
 \qquad \omega>0,\quad\gamma\ge0.
 \label{eq:rr-qubit}
\end{equation}
The depolarizing term has Pauli jump operators
$\sqrt{\gamma/4}\,X$, $\sqrt{\gamma/4}\,Y$, and
$\sqrt{\gamma/4}\,Z$.
Prepare $j=|0\rangle$ or $i=|1\rangle$ and detect
$\Pi_j=|0\rangle\langle0|$.  The instantaneous contrast is
$e^{-\gamma t}\cos(\omega t)$, so uniform timing yields
\begin{equation}
 C(T)=\frac{\gamma+e^{-\gamma T}
       [-\gamma\cos(\omega T)+\omega\sin(\omega T)]}
       {(\gamma^2+\omega^2)T}.
 \label{eq:rr-qubit-contrast}
\end{equation}
For $\gamma=0$, this is $\sin(\omega T)/(\omega T)$.

\begin{theorem}[Orthogonal-preparation damping threshold]
\label{thm:rr-threshold}
For $\gamma>0$, a horizon with $C(T)<0$ exists if and only if
\begin{equation}
 \frac{\omega}{\gamma}\log\frac{\omega}{\gamma}>\frac{3\pi}{2},
 \quad\text{equivalently}\quad
 \frac{\omega}{\gamma}>e^{W(3\pi/2)}
 =3.644173671645632\ldots .
 \label{eq:rr-threshold}
\end{equation}
Here $W$ is the principal real Lambert function.  This threshold concerns
\eqref{eq:rr-qubit}, the stated orthogonal preparations, and the decreasing-clock
class.  It is not a universal coherence threshold.
\end{theorem}

\begin{proof}
Write $I(T)=T C(T)$.  Then $I'(T)=e^{-\gamma T}\cos(\omega T)$.
Its first local minimum following the initial maximum occurs at
$T_*=3\pi/(2\omega)$; subsequent minima occur one full period apart and have
values
\[
 I(T_n)=\frac{\gamma-\omega
 e^{-\gamma(3\pi/2+2\pi n)/\omega}}{\gamma^2+\omega^2},
 \qquad n=0,1,\ldots .
\]
These values increase with $n$.  Since $I(0)=0$, a negative value exists exactly
when $\omega e^{-3\pi\gamma/(2\omega)}>\gamma$.
Setting $r=\omega/\gamma$ gives $r\log r>3\pi/2$, whose unique root above one
is $e^{W(3\pi/2)}$.  The mixture theorem proves that decreasing clocks cannot
produce a reversal when every uniform horizon is nonnegative.
\end{proof}

When reversal exists, the most negative uniform contrast occurs within
$(\pi/\omega,3\pi/(2\omega))$.  Indeed
$C'(T)=[T I'(T)-I(T)]/T^2$, and the numerator is strictly increasing throughout
that interval because its derivative is $T I''(T)>0$.  It changes sign between
the first zero of $I$ and its first negative minimum.  Later periods increase
a negative numerator $I$ toward its positive limiting value while also increasing
$T$, so they cannot improve the negative average.

With $\gamma=1/10$, $\omega=2$, $T=9/4$, and a common affine detector of gain
$9/10$, the observed contrast is
\[
 C_{\rm obs}=-0.1440713137250376\ldots<-\frac{18}{125}.
\]
The strict rational inequality is certified by an $80$-degree Taylor enclosure,
independent of floating-point agreement.  Mean evolution time is $9/8$ per
record.  Preparation, readout, reset, and control energy remain separate costs.
An exponentially distributed readout with rate $\lambda$, by contrast, yields
\[
 C_{\exp}=\frac{\lambda(\lambda+\gamma)}
                 {(\lambda+\gamma)^2+\omega^2}>0.
\]
It respects the classical protection but has no power against this alternative.

\section{The best comparison preparation}
\label{sec:rr-spectral-preparation}

For a quantum channel semigroup $\Phi_t$ and a trusted reference state $|j\rangle$,
define the averaged accessible effect
\[
 A_T=\frac1T\int_0^T\Phi_t^\dagger(\Pi_j)\,dt.
\]
The reference score is $\langle j|A_T|j\rangle$; arbitrary quantum state
preparation attains the maximal score $\lambda_{\max}(A_T)$.  The Rayleigh
variational principle therefore gives the optimal reversal gap
\begin{equation}
 V_T=\lambda_{\max}(A_T)-\langle j|A_T|j\rangle.
 \label{eq:rr-spectral-gap}
\end{equation}
This is established spectral optimization applied to experimental design, not a
new eigenvalue theorem.  A nonzero matrix element connecting $j$ to its
orthogonal subspace guarantees $V_T>0$.

For \eqref{eq:rr-qubit}, write
\begin{align*}
 C&=\frac1T\int_0^T e^{-\gamma t}\cos(\omega t)\,dt,\\
 S&=\frac1T\int_0^T e^{-\gamma t}\sin(\omega t)\,dt
 =\frac{\omega-e^{-\gamma T}
 [\omega\cos(\omega T)+\gamma\sin(\omega T)]}
 {(\gamma^2+\omega^2)T}.
\end{align*}
Then $A_T=(I+C Z+S Y)/2$ for the Hamiltonian convention in
\eqref{eq:rr-qubit}, and
\[
 V_T=\frac{\sqrt{C^2+S^2}-C}{2}.
\]
For $\gamma>0$, the integrated sine $S$ is positive at every $T>0$: successive
positive sine lobes precede their negative partners and carry larger exponential
weight.  Within a negative lobe the smallest accumulated value is its endpoint,
where that paired contribution remains positive.  Thus $V_T>0$ for every
$\gamma,\omega,T>0$ when the comparison state can be coherently optimized.
This does not promise useful precision under arbitrarily strong damping;
for fixed $T$ and $\gamma\gg\omega$, the gap scales approximately as
$\omega^2/(4\gamma^3T)$.

A fixed $+Y$ eigenstate is often nearly optimal and avoids a model-dependent
preparation angle.  Its gap is $(S-C)/2$.  At $\gamma=1/10$, $\omega=2$,
$T=7/4$, and detector gain $9/10$, exact rational enclosures give
\[
 V_{+Y,\rm obs}=0.257271133487835\ldots>\frac{257}{1000}.
\]
Mean evolution time is $7/8$.  Coherent comparison preparation is an additional
physical resource, and already presumes the relevant coherent controls; this
example is not evidence for device-independent quantum certification.

\section{Calibration and established experimental baselines}
\label{sec:rr-comparison}

A common fixed binary detector with response $m(p)=d_0+k p$, $k>0$, preserves
the classical sign and cancels its offset.  Reference-preparation total-variation
error $\epsilon_j$ weakens the classical null by at most $\epsilon_j$.
Comparison-preparation error is unnecessary for that null when any initial
distribution is allowed, but contributes to predicted alternative power.
A common readout-law discrepancy of total variation $\delta$ changes a contrast
by at most $2\delta$.  Separate detector probability errors add their two radii.
A joint calibrated record-law bound should replace, rather than duplicate,
errors already accounted for individually.  Statistical confidence cannot
establish the reset or clock assumptions from the same test records.

A close prior comparator is the survival inequality
\cite{rrLambert2010}
\[
 |2P_{jj}(t)-P_{jj}(2t)|\le1.
\]
The classical semigroup actually gives the stronger relation
\begin{equation}
 P_{jj}(2t)=\sum_kP_{jk}(t)P_{kj}(t)\ge P_{jj}(t)^2.
 \label{eq:rr-survival-product}
\end{equation}
The usual upper inequality follows from $2p-p^2\le1$.  To estimate
\eqref{eq:rr-survival-product} directly, multiply the binary outcomes of two
independent length-$t$ records and compare that product with a length-$2t$
record.  The product is Bernoulli with mean $P_{jj}(t)^2$ and consumes two
physical preparations.

For a linear contrast $\sum_k a_k p_k$ with positive alternative gap $g$, a
common sufficient Hoeffding design has resource proxy
\begin{equation}
 \mathcal B=\frac{\left(\sum_k |a_k|\sqrt{c_k}\right)^2}{g^2},
 \label{eq:rr-hoeffding-proxy}
\end{equation}
where $c_k$ is the resource cost of one independent observation of arm $k$.
The same confidence factor multiplies this proxy for all methods.  This is a
comparison of sufficient concentration-bound designs, not their exact minimum
sample complexities.  For randomized clocks, evolution-time cost below means
\emph{expected} evolution time; it is not a hard bound on elapsed laboratory
time or a measurement of energy.

Table~\ref{tab:rr-baselines} fixes $\gamma=0.1$, $\omega=2$, ideal known readout,
independent records, and zero preparation overhead.  Each column optimizes its
own time parameter numerically on $[0.001,5]$.  Physical record costs for the
product arm include both preparations.  The coherent-preparation variants use
an additional control resource.

\begin{table}[htbp]
\centering
\begin{tabular}{lrr}
\hline
Method & Record proxy & Expected-time proxy\\
\hline
Uniform clock, orthogonal preparation &156.08&174.02\\
Established linear survival inequality &219.68&133.64\\
Classical survival-product inequality &75.84&58.02\\
Uniform clock, fixed quadrature preparation &48.81&40.15\\
Uniform clock, optimized preparation &48.39&40.02\\
\hline
\end{tabular}
\caption{Matched sufficient-bound comparisons.  Lower is better.  These are
numerical design results, not physical measurements or certified global
statistical optima.  The survival-product baseline beats the simplest new
assay; extra coherent preparation changes the available operations.}
\label{tab:rr-baselines}
\end{table}

The comparison rejects a claim of universal superiority for the randomized
assay.  Its contribution is a different exact access contract, protected timing
families, and a route to calibrated design.  Which method is experimentally
preferable depends on preparation access, detector errors, timing, and actual
apparatus costs.  A framework becomes more useful when it selects among these
methods and records their failure cases, not when it relabels every modeled
quantum signal as a breakthrough.

\section{Preparation design for a baseline-subtracted dimension witness}
\label{sec:rr-baseline-preparation}

The dimension result developed later, Proposition~\ref{rr:dimension-robust},
changes the null assumptions again.  For a homogeneous CTMC with at most three
hidden states, it permits arbitrary fixed preparations $\mu,\nu$ and an
arbitrary fixed stochastic detector $f\in[0,1]$.  Writing
\[
 C_0=(\mu-\nu)f,\qquad
 C_T=(\mu-\nu)\frac1T\int_0^T e^{tQ}f\,dt,
\]
its bound is $|C_0-C_T|\le1$.  Separately measuring the two preparations at
baseline and at randomized time therefore removes the exact hidden-reset
assumption of the earlier assay, at the price of a restriction on classical
hidden-state dimension.  Section~\ref{rr:four-states} exhibits the importance
of that restriction.  The quantum design calculation is independent of the
proof of the classical bound.

\begin{lemma}[Quantum preparation design for the dimension witness]
\label{lem:rr-dimension-preparation}
For a quantum averaged effect $A_T$ and the same baseline effect $\Pi_j$, let
$D_T=\Pi_j-A_T$.  The largest value of
\[
 W=\operatorname{Tr}\!\left[D_T(\rho_+-\rho_-)\right]
\]
over two arbitrary quantum preparations is the spectral diameter
$\lambda_{\max}(D_T)-\lambda_{\min}(D_T)$.
For the qubit model \eqref{eq:rr-qubit}, it is
\begin{equation}
 W_{\max}(T)=\sqrt{(1-C)^2+S^2},
 \label{eq:rr-dimension-optimum}
\end{equation}
attained by antipodal pure Bloch vectors parallel to $(0,-S,1-C)$.
\end{lemma}

\begin{proof}
The maximum and minimum expectations of a Hermitian matrix are its largest and
smallest eigenvalues, independently attainable by the two preparations.
Here $D_T=[(1-C)Z-SY]/2$, whose eigenvalues are
$\pm\sqrt{(1-C)^2+S^2}/2$.
\end{proof}

This is a standard spectral variational calculation used for experimental
selection.  Orthogonal $Z$ preparations instead give $W=1-C$, whose optimal
horizon is exactly the most negative uniform contrast studied above.  A common
affine detector of gain $k$ multiplies these modeled quantum contrasts by $k$.
For those orthogonal preparations, a predicted violation requires
$k>1/(1-C)$.  The classical dimension bound itself does \emph{not} require this
affine detector model or a known positive gain: it already permits any fixed
stochastic $f$.  Gain enters prediction of experimental power, not validity
of that null.

A simple rational preparation gives a larger certified margin than the
orthogonal pair.  Set
\[
 \gamma=\frac1{10},\quad\omega=2,\quad T=2,\quad k=\frac9{10},
 \qquad n=\left(0,-\frac{12}{37},\frac{35}{37}\right).
\]
The vectors $n$ and $-n$ are pure because $12^2+35^2=37^2$.  Exact rational
Taylor enclosures certify
\begin{equation}
 W=k\left[\frac{35}{37}(1-C)+\frac{12}{37}S\right]
 =1.0806050936582667\ldots>\frac{27}{25}.
 \label{eq:rr-rational-dimension-witness}
\end{equation}
Each randomized-time arm has mean evolution time $1$; each baseline arm has
zero evolution time in the ideal timing convention.  Four equally sampled
arms therefore have mean evolution time $1/2$ per physical record.  All four
arms still incur preparation, measurement, and reset costs.  A common
uncontrolled readout delay must be included consistently in the baseline and
random-time effects; it is not automatically covered by the zero-delay
alternative prediction above.

Numerical optimization gives an ideal-state value about $1.20123$ at
$T\simeq2.04524$ for this damping and frequency, or $1.08111$ after gain $0.9$.
The rational design is close to that numerical value, but the reported search
is not a certificate of a global optimum.  Established dimension witnesses and
survival tests remain relevant comparators.  The present result is an explicit
physical alternative and a reproducible preparation choice for the stated
finite-dimension null, not an unrestricted exclusion of classical hidden
mechanisms.
